\documentclass[11pt,a4paper]{altacv}

\geometry{left=1.5cm,right=1.5cm,top=1.cm,bottom=1.2cm,columnsep=1.5cm}

\usepackage[utf8]{inputenc}
\usepackage[T1]{fontenc}
\usepackage[french]{babel}
\usepackage{paracol}
\usepackage{xcolor}
\usepackage{hyperref}
\usepackage{fontawesome5}
\usepackage{setspace}

\definecolor{SlateGrey}{HTML}{2E2E2E}
\definecolor{LightGrey}{HTML}{666666}
\definecolor{DarkPastelRed}{HTML}{8AC0D9}
\definecolor{PastelRed}{HTML}{00B0DA}
\definecolor{GoldenEarth}{HTML}{B4AD0C}
\colorlet{name}{black}
\colorlet{tagline}{PastelRed}
\colorlet{heading}{DarkPastelRed}
\colorlet{headingrule}{GoldenEarth}
\colorlet{subheading}{PastelRed}
\colorlet{accent}{PastelRed}
\colorlet{emphasis}{SlateGrey}
\colorlet{body}{LightGrey}

\renewcommand{\familydefault}{\sfdefault}

\name{\Large Guillaume BOILEAU}
\tagline{Chef de projet technique | Validation logicielle, IVVQ, spatial et systèmes complexes}
\personalinfo{
  \email{guillaume.boileaupro@gmail.com}
  \phone{+33 6 59 94 85 84}
  \location{Alpes-Maritimes, France}
  \linkedin{guillaume-boileau-phd-891b81265}
  \researchgate{Guillaume-Boileau-2}
  \orcid{0000-0002-3576-6968}
}

\setlength{\parskip}{0.75\baselineskip}

\begin{document}
\makecvheader
\vspace{-0.6cm}

\cvsection{Lettre de motivation}
\vspace{-0.4cm}

{\color{accent}\textbf{Objet :}} \textit{Candidature au poste de Manager IVVQ Logiciel - Secteur spatial H/F}

{\color{body}

Madame, Monsieur,

Je vous adresse ma candidature pour le poste de Manager IVVQ Logiciel - Secteur spatial au sein de Viveris. Cette opportunité m'intéresse particulièrement car elle se situe au croisement de plusieurs dimensions qui structurent mon parcours : la validation de logiciels complexes, l'intégration de chaînes techniques, l'analyse d'anomalies, la traçabilité des résultats, la qualité des livrables, le reporting technique et la coordination entre équipes logiciel, système et contributeurs techniques.

Viveris intervient sur des projets d'ingénierie exigeants, notamment dans le domaine des systèmes embarqués critiques et du spatial. C'est précisément ce type d'environnement qui m'attire : des systèmes à forte contrainte, des interactions nombreuses entre composants, une nécessité de rigueur dans les procédures, et une responsabilité forte sur la qualité des résultats produits. Votre besoin d'un profil capable à la fois de contribuer aux activités opérationnelles IVV et de piloter les activités de validation correspond bien à la manière dont mon parcours s'est construit.

Mon expérience s'est développée à l'interface entre systèmes complexes, développement logiciel, analyse de données, validation, coordination technique et pilotage de contributions. J'ai travaillé dans des contextes où la fiabilité des outils, la reproductibilité des résultats, la documentation, la structuration des workflows et la coordination entre plusieurs acteurs sont essentielles. Cette culture de travail m'a appris à ne pas seulement produire du code ou des analyses, mais aussi à vérifier leur cohérence, organiser leur validation, documenter les résultats et assurer leur exploitation par d'autres contributeurs.

Au Laboratoire Lagrange de l'Observatoire de la Côte d'Azur, dans le cadre de travaux menés avec le CNES, j'ai contribué à des activités de développement, d'intégration et de validation pour la mission spatiale LISA, mission ESA/NASA de grande envergure. Dans ce cadre, j'ai notamment développé CATpyLISA, une plateforme Python destinée à l'intégration, la comparaison et la validation de modèles et de chaînes d'analyse complexes. Ce travail m'a amené à structurer des workflows techniques, intégrer des composants hétérogènes, organiser des campagnes de comparaison, analyser des incohérences, documenter les résultats et coordonner des contributions dans un cadre international.

Cette expérience est directement transposable aux activités IVVQ décrites dans votre offre. Elle m'a permis de définir et d'exécuter des démarches de validation, de vérifier la cohérence de sorties logicielles complexes, d'analyser des comportements inattendus, de consolider les résultats, de produire des éléments de reporting et de maintenir une documentation exploitable. J'ai également développé une forte sensibilité à la traçabilité : il ne suffit pas d'obtenir un résultat, il faut pouvoir expliquer son origine, documenter les hypothèses, identifier les écarts, suivre les corrections et garantir que les livrables restent compréhensibles et réutilisables dans la durée.

Dans le cadre de LISA, j'ai aussi été confronté à des chaînes logicielles et scientifiques complexes, dans lesquelles les résultats dépendent de nombreux paramètres, modèles, outils et procédures. Cette situation impose une approche proche de l'IVVQ : définir ce qui doit être vérifié, construire des scénarios de comparaison, identifier les écarts, comprendre les causes possibles, suivre les corrections et documenter les conclusions. Même si le contexte n'était pas celui d'un logiciel embarqué satellite au sens strict, les méthodes mobilisées - intégration, validation, analyse d'anomalies, automatisation Python, contrôle de cohérence, documentation et coordination - sont directement pertinentes pour des activités de validation logicielle exigeantes.

J'ai également encadré plusieurs stagiaires de niveau Master et coordonné des activités impliquant ingénieurs, techniciens et chercheurs. Cette dimension est importante pour le poste proposé, qui demande non seulement de contribuer techniquement, mais aussi de piloter et accompagner une équipe IVV. J'ai appris à clarifier les objectifs, répartir les actions, suivre l'avancement, aider les contributeurs à monter en autonomie, maintenir l'alignement technique et produire des points de synthèse utiles à la prise de décision. Cette expérience m'a donné une posture de coordination opérationnelle : rester proche du contenu technique tout en gardant une vision d'ensemble sur les priorités, les risques et les livrables.

Plus récemment, chez VEV Platform Services France, j'ai travaillé sur des services logiciels liés aux infrastructures de recharge électrique, dans un contexte industriel connecté à des systèmes terrain. J'y ai contribué au développement de services web et backend en TypeScript, Angular et Node.js, ainsi qu'au suivi de flux de données opérationnels et d'interactions avec des équipements matériels. Cette expérience m'a permis de renforcer ma compréhension des contraintes liées aux systèmes réels : incidents en production, échanges de données, fiabilité opérationnelle, interfaces entre plateforme logicielle et matériel, et nécessité de résoudre les problèmes avec méthode.

Ce parcours me permet aujourd'hui d'apporter une double culture utile pour Viveris : une culture spatiale et systèmes complexes acquise dans le cadre de LISA, et une culture logicielle industrielle renforcée par mon expérience récente sur des services opérationnels. Je suis à l'aise dans les environnements où il faut comprendre rapidement une architecture, identifier les points critiques, structurer une démarche de validation, suivre les anomalies, documenter les résultats et travailler avec plusieurs interlocuteurs techniques.

Je souhaite être transparent sur mon positionnement : je ne me présente pas comme un expert bas niveau des architectures ARMv8, de Yocto, des boot softwares ou des drivers embarqués. En revanche, je possède une expérience solide en validation logicielle, intégration de chaînes complexes, automatisation Python, environnements Linux/Unix, scripting, suivi d'anomalies, documentation, reporting, coordination technique et compréhension des interactions entre logiciel, système et matériel. Ce sont des compétences que je peux mobiliser rapidement pour contribuer aux activités IVV, au pilotage de campagnes de validation, au suivi des non-régressions, à la qualité des livrables et à l'organisation du travail d'équipe.

Le poste proposé par Viveris m'intéresse également parce qu'il comporte une dimension de garantie de la qualité des livrables. C'est un aspect auquel je suis particulièrement attaché. Dans mes expériences précédentes, j'ai souvent travaillé sur des sujets où un résultat mal documenté, mal tracé ou difficile à reproduire perd rapidement de sa valeur. J'ai donc développé une manière de travailler fondée sur la clarté des procédures, la structuration des analyses, la consolidation des résultats et la transmission d'informations exploitables par les autres membres du projet.

Je pourrais contribuer à vos projets en apportant une capacité à structurer les activités de validation, organiser les priorités, formaliser les points ouverts, analyser les anomalies, automatiser certaines vérifications, produire une documentation claire et faciliter la coordination entre les équipes logiciel, système et projet. Mon objectif serait de m'inscrire dans une dynamique d'équipe, en apportant à la fois rigueur technique, capacité d'analyse et sens de la coordination.

Rigoureux, autonome et habitué aux environnements internationaux, je suis motivé par les projets où la qualité technique, la validation et la coordination jouent un rôle central. Rejoindre Viveris représenterait pour moi l'opportunité de mettre mon expérience au service de systèmes embarqués critiques dans le secteur spatial, tout en continuant à évoluer dans un environnement d'ingénierie exigeant, collaboratif et orienté résultats.

Je serais heureux de pouvoir échanger avec vous afin de vous présenter plus en détail mon parcours, mon positionnement et la manière dont je pourrais contribuer aux activités IVVQ de vos projets spatiaux.

Je vous prie d'agréer, Madame, Monsieur, l'expression de mes salutations distinguées.

\textbf{Guillaume Boileau}

}

\end{document}